\chapter{Preparing Meshes}
\label{chap:mpas_grid_preparation}

This chapter describes the steps used to prepare SCVT meshes for use in MPAS-Atmosphere.
For quasi-uniform meshes, very little preparation is actually needed, and
generally, one only needs to prepare mesh decomposition files --- files that
describe the decomposition of the SCVT mesh across processors --- when running
MPAS-Atmosphere using multiple MPI tasks. The procedure for creating these mesh
decomposition files is described in the first section. 

For variable-resolution SCVT meshes, the area of mesh refinement may be rotated
to any part of the sphere using a program, grid\_rotate, described in the second
section. This utility program may be obtained from the MPAS-Atmosphere download page.

When running limited-area simulations, the limited-area mesh may be produced by
subsetting the elements in an existing SCVT mesh as described in the third section.

\section{Graph partitioning}
Before MPAS can be run in parallel, a mesh decomposition file with an
appropriate number of partitions (equal to the number of MPI tasks that will be
used) is required. A limited number of mesh decomposition files, named {\tt
graph.info.part.*}, are provided with each mesh, as is the mesh
connectivity file, named {\tt graph.info}. If the number of MPI tasks to be used when
running MPAS matches one of the pre-computed decomposition files, then there
is no need to run METIS.
\subsection{Pre-computing graph partition files with METIS} 
\label{sec:metis}

In order to create new mesh decomposition files for some particular number of MPI tasks, 
only the {\tt graph.info} file is required.  The currently supported method for partitioning
a graph.info file uses the METIS software
(\url{http://glaros.dtc.umn.edu/gkhome/views/metis}).  The serial graph partitioning
program, METIS (rather than ParMETIS or hMETIS) should be sufficient for
quickly partitioning any mesh usable by MPAS.

After installing METIS, a {\tt graph.info} file may be partitioned into $N$
partitions by running

\vspace{12pt}
{\tt > gpmetis -minconn -contig -niter=200 graph.info} $N$
\vspace{12pt}

\noindent where $N$ is the required number of partitions. The resulting file, {\tt graph.info.part.}$N$, 
can then be copied into the MPAS run directory before running the model with $N$ MPI tasks.

\subsection{Online graph partitioning with Scotch}
\label{sec:scotch}
Starting with the v8.4.0 release, MPAS also provides an option to enable online graph-partitioning feature using the Scotch library
(\url{https://www.labri.fr/perso/pelegrin/scotch/}). To be able to use this feature, the Scotch library
must first be installed on the system, and then MPAS must be built with Scotch support. The minimum supported SCOTCH version is v7.0.8.

If these conditions are met, then the model can partition the mesh at runtime, using only the connectivity information
present in the input file, and without the need for any {\tt graph.info} files. After computing the graph partitioning,
MPAS will write the generated partition file to disk for future use.

In practical terms, online partitioning is usually most beneficial when starting from a fresh run directory, changing the
MPI task count frequently, or running workflows where pre-generating many {\tt graph.info.part.*} files is inconvenient.

Internally, PT-Scotch performs this operation through distributed graph mapping algorithms. In this framework,
partitioning into $N$ parts is treated as mapping the mesh graph onto a target graph with $N$ domains. This
approach allows MPAS to build partitions in parallel, while still enforcing load-balance and connectivity goals.
 
The exact behavior of the online graph partitioning can be controlled through namelist options.

\subsubsection{Building Scotch}
To build Scotch, first download the source code from the Scotch website. You can either download and extract release
tarballs from this link:
\url{https://gitlab.inria.fr/scotch/scotch/-/releases}, or clone the git repository using the command below:

\vspace{12pt}
{\tt > git clone https://gitlab.inria.fr/scotch/scotch.git \&\& cd scotch} \\
{\tt > git checkout v7.0.8}
\vspace{12pt}

To be able to use the distributed graph partitioning provided by PT-Scotch, you will need to pass an appropriate MPI installation path during the Scotch build.

Other system pre-requisites are Bison and Flex. Note that a more recent version of Flex (>v2.6.4) is recommended. Scotch documentation also references the requirement for a Bison version > 3.4.

\subsubsection{Building MPAS with Scotch}
After building and installing Scotch, with a minimum version of at least 7.0.8, as described above, the MPAS atmosphere and init\_atmosphere cores can be built with online graph partitioning support by setting the path to Scotch prior to building MPAS.

\vspace{12pt}
{\tt > export SCOTCH=/path/to/scotch/installation} \\
{\tt > make nvhpc CORE=atmosphere}
\vspace{12pt}

If MPAS is built and linked with Scotch successfully, a message indicating Scotch support will appear near the
end of the build output.

\subsubsection{Usage and run-time behavior}
After building MPAS with Scotch, the user can still choose between using a pre-generated graph partition file or the SCOTCH online graph partitioning, by setting appropriate values for the {\tt config\_block\_decomp\_file\_prefix} namelist option.

If {\tt config\_block\_decomp\_file\_prefix+mpi\_tasks} points to a valid graph partition file that already exists in the run directory, then it is used to proceed with the model run without invoking Scotch.

If {\tt config\_block\_decomp\_file\_prefix==''} or {\tt config\_block\_decomp\_file\_prefix+mpi\_tasks} does not match any valid graph partition file in the run directory, then Scotch graph partitioning is invoked. During this process, the generated partition is saved as a graph partition file to the run directory so that it may be reused in the following runs without needing to invoke Scotch again.

These run-time choices are summarized in Table~\ref{tab:scotch_runtime_behavior}.

If {\tt config\_block\_decomp\_file\_prefix+mpi\_tasks} points to a valid graph partition file that already exists in the run directory, then it is used to proceed with the model run without invoking Scotch.

If {\tt config\_block\_decomp\_file\_prefix==''} or {\tt config\_block\_decomp\_file\_prefix+mpi\_tasks} does not match any valid graph partition file in the run directory, then Scotch graph partitioning is invoked. During this process, the generated partition is saved as a graph partition file to the run directory so that it may be reused in the following runs without needing to invoke Scotch again. These run-time choices are summarized in Table~\ref{tab:scotch_runtime_behavior}.

\vspace{12pt}
\begin{longtable}{|p{2.9in}|p{2.9in}|}
\hline
\textbf{Condition} & \textbf{Run-time behavior} \\ \hline
{\tt config\_block\_decomp\_file\_prefix+mpi\_tasks} points to a valid graph partition file that already exists in the run directory &
The existing graph partition file is used to proceed with the model run without invoking Scotch. \\ \hline
{\tt config\_block\_decomp\_file\_prefix==''} or {\tt config\_block\_decomp\_file\_prefix+mpi\_tasks} does not match any valid graph partition file in the run directory &
Scotch graph partitioning is invoked. The generated partition is saved to the run directory for reuse in later runs. \\ \hline
\caption{SCOTCH online partitioning run-time behavior}\label{tab:scotch_runtime_behavior} \\
\end{longtable}
\vspace{12pt}

If MPAS has not been built with SCOTCH, any code paths relating to Scotch will not be taken. And an incorrect specification of {\tt config\_block\_decomp\_file\_prefix+mpi\_tasks} should lead to the model halting.


\section{Relocating refinement regions on the sphere}
\label{sec:grid_rotate} 

The purpose of the grid\_rotate program is simply to rotate an MPAS mesh file,
moving a refinement region from one geographic location to another, so that the
mesh can be re-used for different applications. This utility was developed out
of the need to save computational resources, since generating an SCVT ---
particularly one with a large number of generating points or a high degree of
refinement --- can take considerable time.

To build the grid\_rotate program, the Makefile should first be edited to set
the Fortran compiler to be used; if the netCDF installation pointed to by the
{\tt NETCDF} environment variable was build with a separate Fortran interface
library, it will also be necessary to add {\tt -lnetcdff} just before {\tt -lnetcdf} in 
the Makefile. After editing the Makefile, running `make' should
result in a grid\_rotate executable file.

Besides the MPAS grid file to be rotated, grid\_rotate requires a namelist file,
{\tt namelist.input}, which specifies the rotation to be applied to the mesh.
The namelist variables are summarized in the table below
   
\vspace{12pt}
\begin{longtable}{|p{3.25in} |p{2.5in}|}
\hline
config\_original\_latitude\_degrees & original latitude of any point on the sphere \\ \hline
config\_original\_longitude\_degrees & original longitude of any point on the sphere \\ \hline
config\_new\_latitude\_degrees &  latitude to which the original point should be shifted \\ \hline
config\_new\_longitude\_degrees &  longitude to which the original point should be shifted \\ \hline
config\_birdseye\_rotation\_counter\_clockwise\_degrees & rotation about a vector from the sphere center through the original point \\ \hline
\end{longtable}
\vspace{12pt}

\noindent Essentially, one chooses any point on the sphere, decides where that
point should be shifted to, and specifies any change to the orientation (i.e.,
rotation) of the mesh about that point. 

Having set the rotation parameters in the {\tt namelist.input} file, the
grid\_rotate program should be run with two command-line options specifying the
original grid file name and the name of the rotated grid file to be produced,
e.g.,

\vspace{12pt}
{\tt > grid\_rotate grid.nc grid\_SE\_Asia\_refinement.nc}
\vspace{12pt}

The original grid file will not be altered, and a new, rotated grid file will be
created. The NCL script {\tt mesh.ncl} may be used to plot either of the
original or rotated grid files after suitable setting the name of the grid file
in the script.

\vspace{12pt}
{\em Note: The grid\_rotate program initializes the new, rotated grid file to a copy of the original grid file.
If the original grid file has only read permission (i.e., no write permission), then so will the copy, and
consequently, the grid\_rotate program will fail when attempting to update the fields in the copy.}


\section{Creating limited-area SCVT meshes}
\label{sec:mesh_subset} 

The process of creating a limited-area (regional) mesh for MPAS-Atmosphere involves
selecting any existing mesh -- either a quasi-uniform mesh, or a variable-resolution mesh
that has been rotated as in section \ref{sec:grid_rotate} -- describing the geographical
region to be extracted from that mesh, and running the limited-area Python program to extract
all cells, edges, and vertices in the designated region. The result is a new netCDF mesh file that can
be used to make a limited-area simulation as described in Section \ref{sec:regional}.

The limited-area Python program may be obtained from the MPAS-Atmosphere download page. Although
the set of required Python packages may change over time, the program currently requires
the {\tt numpy} and {\tt netCDF4} packages, in addition to other standard packages.
